% Sección 5. Fuentes: documentos/informe_resultados.md (sección 8, «Limitaciones y preguntas abiertas»),
% app/DATOS.md (contradicciones C8, C9 y C17) y los README de E7, D5, D6, E12, E13 y sistema_tiempo_real/.
\section{Discusión y limitaciones}
\label{sec:discusion}

El pronóstico de corto plazo, que es el del enunciado, quita el \CPRutaQuitaLX~\% del error del promedio por día y
hora por ruta a 60 minutos, en meses que no se usaron para elegir nada (sección~\ref{sec:corto-plazo}). El
pronóstico semanal mejora a la regla manual en unos seis puntos de WAPE, con una ventaja que resiste la antelación, la
ruta y el año. Pero cuántos buses despachar lo decide lo que cuesta quitar una salida: con $\beta \approx 0{,}70$,
con la conservadora no paga ningún plan y con la simétrica pagan los dos; recortar deja más caja a costa de pasajeros y
sobrecarga. Las limitaciones principales:
\begin{itemize}
  \item \textbf{$\beta$ viene de fallas imprevistas y pequeñas} (un 5~\% de las salidas de la hora en promedio). Un
        recorte planificado y parejo podría tener una $\beta$ de 0,4--0,5, y entonces recortar sí pagaría; si los viajes
        \comi{sin registro} corrieron, $\beta$ estaría más cerca de 1. Los pasajeros ganados al añadir salidas no se han
        medido: la contabilidad conservadora es un piso y la simétrica un techo.
  \item \textbf{Costos con fuente, pero no de Sobusa}: rendimiento, mantenimiento y pago de los conductores salen de
        estudios de otras ciudades, y la espera y la penalidad por sobrecarga son supuestos. Los conductores regulares
        registran unas 247~h en ruta al mes, por encima de las 182~h de la jornada ordinaria de 2026.
  \item \textbf{Optimizador y simulador}: el optimizador se evalúa con su propia función y con la demanda servida bajo la
        oferta real, y no ve el hacinamiento por bus; el simulador llena los buses más que la realidad y no tiene
        elasticidad.
  \item \textbf{Selección y objetivo}: elegir entre 51 candidatos en los mismos meses infla algo la ventaja del ganador
        (Holm y la prueba en 2024 lo acotan; el pronóstico de la sección~\ref{sec:corto-plazo} no tiene ese
        problema, porque se eligió y se reportó con meses distintos), y los pronósticos de E1 a E12, y el plan del
        motor, usan el objetivo
        viejo, que subestima lo que sube a los buses entre un 5,6 y un 7,7~\%.
  \item \textbf{Datos y tiempo real}: 2026 está sin confirmar; once de los doce últimos días de mes de 2025 siguen sin
        eventos de conteo; no se conoce la latencia de los datos de Sobusa, y la corrida vigente del backtest y el modelo
        persona a persona no tuvieron revisor independiente (sí la corrida anterior, con el mismo código del motor).
  \item \textbf{Pronóstico de corto plazo}: en los domingos corrientes no mejora al promedio por día y hora, sobre
        todo en CRA~54; el sentido es una aproximación (la mitad del ciclo) y la regla del
        segmento más cercano deja parte de las subidas en el sentido contrario; la prueba de julio a diciembre de
        2025 no es ciega en sentido estricto, porque esos meses se conocían por los experimentos anteriores, con
        otro objetivo; el modelo supone que el conteo llega en vivo, y la API, que ya lo sirve, lee todavía el
        panel de 2023 a 2025 y no la captura en vivo.
\end{itemize}
